%%%%%%%%%%%%%%%%%%%%
%
% $Autor: Wings $
% $Datum: 2020-02-24 14:29:03Z $
% $Pfad: komponenten/Bilderkennung/Produktspezifikation/CorelTPU/Ausarbeitung/Kapitel/Software.tex $
% $Version: 1791 $
%
%
%%%%%%%%%%%%%%%%%%%




%todo neu formulieren
\subsection{Nutzung der Kamera in Python}

%todo
Wie wird die Kamera unter Python verwendet?



\subsection{Zeitverhalten}

%todo Überabreiten; Quellen 
Die SmartPiCam benötigt für einen Durchlauf der Schleife von Schritt 6 bis 15 170–180ms. Das entspricht knapp 6\ac{fps}. Der Mensch nimmt Objekte zum Vergleich laut Wikipedia in 200–300ms wahr. 35–40ms werden für die Anzeige der Fotos benötigt und circa 60–75ms entfallen auf das Inferencing auf dem Edge Computer inklusive Datenübertragung (rund 264KB). Das entspricht einer Auslastung des Edge TPU  von gerade einmal 33–42\%. Entfernt man die Anzeige der Fotos aus dem Programm, so steigt die Bildrate weiter auf 7,3fps. Diese Programmvariante 
heißt \FILE{smartPiCamMainNoDisp.py} und ist auch im 
GitHub-Repository zum Projekt enthalten. Die Edge TPU 
hat hier also noch große Reserven, wenn man bedenkt, 
dass sie auch auf doppelte Geschwindigkeit konfiguriert 
werden kann. Letzteres macht aber erst richtig Sinn, wenn 
ein Programm auch Fotos in entsprechender Geschwindigkeit liefert.
